\section{Results}
\label{sec:results}

\subsection{Controlled Evidence-Contract Results}

Table~\ref{tab:mutation-results} reports every frozen candidate. All 20 valid
reports passed deterministic verification. All 16 mutated reports were
intercepted, yielding 20/20 valid-report acceptance, 16/16 hard-defect recall,
and 0/20 false blocks on this controlled benchmark. The verifier routed
repairable selector and grounding defects to \textsc{needs-evidence}; it routed
scope, hash-integrity, and diagnosis-conflict defects directly to
\textsc{review-required}.

\begin{table}[t]
    \caption{Deterministic verification on the 36-candidate frozen benchmark.
    ``NE'' denotes \textsc{needs-evidence}; ``HR'' denotes
    \textsc{review-required}.}
    \label{tab:mutation-results}
    \centering
    \small
    \begin{tabular}{@{}lrrr@{}}
        \toprule
        Candidate family & $n$ & NE & HR \\
        \midrule
        Valid, unmodified & 20 & 0 & 0 \\
        Unsupported scope & 6 & 0 & 6 \\
        Invalid selector & 2 & 2 & 0 \\
        Value/hash mismatch & 2 & 0 & 2 \\
        Claim not grounded & 2 & 2 & 0 \\
        Critical span not grounded & 2 & 2 & 0 \\
        Diagnosis conflict & 2 & 0 & 2 \\
        \midrule
        All injected defects & 16 & 6 & 10 \\
        \bottomrule
    \end{tabular}
\end{table}

These results test exact implemented rules against programmed defects. They do
not estimate performance on naturally occurring diagnoses and do not measure
the optional semantic verifier. Their contribution is narrower: the
claim--evidence contract creates an executable boundary that distinguishes
valid reports from six classes of known invalid report.

\subsection{System Validation Results}

Table~\ref{tab:system-results} summarizes the system-level evidence. The local
Windows candidate verification for v0.7.0 collected 1,807 tests: 1,803 passed,
one was skipped, and three failed only because CRLF checkout changed
byte-for-byte frozen reference assets. Statement coverage was 92.14\%, above
the 91\% project threshold, but the three known failures mean this candidate is
not a fully green release gate. Static analysis and packaging were checked
separately. The accepted recovery gate also completed all 20 independent SQLite
stability processes and preserved an end-to-end review decision across an
API-container restart.

\begin{table}[t]
    \caption{SpanVouch engineering validation. Candidate-suite failures are
    limited to the documented Windows CRLF frozen-reference mismatch.}
    \label{tab:system-results}
    \centering
    \small
    \begin{tabular}{@{}p{0.30\columnwidth}p{0.62\columnwidth}@{}}
        \toprule
        Validation surface & Observed result \\
        \midrule
        Code snapshot & Local v0.7.0 candidate on Windows \\
        Public contracts & 6 versioned roots with strict valid/invalid fixtures \\
        Release suite & 1,803 passed, 1 skipped, 3 CRLF-only failures; 92.14\% coverage \\
        Static/delivery gates & Ruff, strict mypy, locked wheel, and Compose passed \\
        Process recovery & 20/20 SQLite processes completed \\
        Container recovery & Unprivileged API retained the terminal decision after restart \\
        Frozen artifact & Manifest and metrics hashes reproduced exactly \\
        \bottomrule
    \end{tabular}
\end{table}

\subsection{Offline Matrix Results}

The offline matrix evaluated all 24 scheduled domain--framework--condition
cells. It completed four adapter executions and both framework-parity
comparisons, then generated the evaluated-results manifest, statistical assets,
claim-gate report, and paper assets. Repeated executions produced identical
bundles. The run attempted zero provider requests and zero GPU operations.

The matrix result establishes that every B0--B5 delivery and analysis interface
is executable across both labs and framework adapters. It does not rank the six
conditions: fake semantic-provider outputs intentionally make the recorded
condition risks unsuitable for an effectiveness comparison.

\subsection{Formal DeepSeek-only Observations}

The v0.7 aggregate bundle is bound to evaluated-results SHA-256
\nolinkurl{bc09f1b134de9370a3b5209fa5e959bce01abbcdf05c8456af1f069fc4cd3088}.
It contains 2,148 scheduled and evaluated plans, with zero missing cells. B0,
B1, B2, and B3 completed; B4 and B5 were \textsc{policy-skipped}, so no Qwen
results are available.

On SupportLab, the paired B3-minus-B2 risk difference was $-0.102391$; its 95\%
interval ran from $-0.165462$ to $-0.046525$. The coverage difference was
$-0.254190$, with limits $-0.358543$ and $-0.155989$. On OpsLab, the preliminary
paired risk difference was $-0.116192$, with limits $-0.206115$ and $-0.032821$.
These are observed DeepSeek-only risk--coverage
tradeoffs under the frozen matrix. They do not establish a causal effect,
cross-model generalization, verifier independence, or an empirical target-risk
certificate. The machine-readable H1--H5 gates remain unresolved.
